\section{Типичные ошибки и их устранение}

В процессе работы с дроном часто возникают ситуации, когда <<код правильный, но ничего не работает>>. В этой главе собраны самые частые проблемы.

\subsection{Проблемы с камерой и маркерами}

\subsubsection{Дрон не видит маркер}
\begin{itemize}
    \item \textbf{Причина 1: Плохое освещение.} Камере темно, выдержка увеличивается, появляется <<смаз>> (Motion Blur) при движении.
    \item \textbf{Решение:} Включите дополнительный свет. Маркер должен быть освещен равномерно.
    
    \item \textbf{Причина 2: Блики.} Маркер заламинирован глянцевой пленкой или лежит под лампой.
    \item \textbf{Решение:} Переместите маркер или используйте матовый скотч.
    
    \item \textbf{Причина 3: Неверный размер.} В коде указан размер 0.2 метра, а напечатан маркер 18 см.
    \item \textbf{Решение:} Измерьте черную рамку линейкой и исправьте константу в коде.
    
    \item \textbf{Причина 4: Quiet Zone.} Маркер вырезан по черному контуру без белого поля.
    \item \textbf{Решение:} Наклейте его на белый лист бумаги.
\end{itemize}

\subsubsection{Дрон летит не туда}
\begin{itemize}
    \item \textbf{Причина 1: Перепутаны оси.} Вы думаете, что X — это вправо, а это вперед.
    \item \textbf{Решение:} Проверьте ориентацию (Body Frame). Ось X дрона смотрит вперед --- туда же, куда нос и передняя камера; Y --- влево, Z --- вверх. Не путайте с системой самой камеры в OpenCV: там вперед <<из объектива>> смотрит ось Z (вектор \texttt{tvec}). И с командами \texttt{...\_body\_fixed} в \texttt{pioneer\_sdk}: в них x --- вправо, y --- вперёд.

    \item \textbf{Причина 1а: Перепутан знак поправки.} Дрон отлетает от метки вместо того, чтобы приближаться к ней.
    \item \textbf{Решение:} Поменяйте знак коэффициента по этой оси и проверьте опытом: после короткого импульса скорости ошибка должна уменьшиться (миссия Python №\,92).

    \item \textbf{Причина 2: Ошибка калибровки.} Камера <<рыбий глаз>> сильно искажает картинку по краям.
    \item \textbf{Решение:} Проведите калибровку камеры заново.
    
    \item \textbf{Причина 3: Yaw (Рысканье).} Дрон повернулся на 90 градусов, но код продолжает отправлять команды в старой системе координат.
    \item \textbf{Решение:} Используйте глобальную систему координат (World Frame) или учитывайте текущий угол Yaw.
\end{itemize}

\begin{taskbox}[Неверный знак]
\emph{Симулятор.} В контроллере \texttt{align} из \verb|staged_landing_sim.py| поменяйте знак поправки по Y (\verb|vy = clamp(KP * ey / h, V_MAX)|) и запустите программу над меткой 12. Опишите, что видно в журнале, и верните знак.
\expected{с неверным знаком $|e_y|$ растёт от кадра к кадру, пока метка не уходит из кадра; контроллер обнуляет скорость, за 120 попыток не достигает центра, и программа печатает «Прервано» и садится, а не улетает дальше.}
\mission{Python}{92, 96}
\end{taskbox}

\subsection{Проблемы с подключением и Wi-Fi}

\subsubsection{Connection Refused / Timeout}
\begin{itemize}
    \item \textbf{Причина 1:} Компьютер подключился к другой Wi-Fi сети (с интернетом).
    \item \textbf{Решение:} Проверьте, что вы подключены к точке доступа Pioneer\_XXXX.
    
    \item \textbf{Причина 2:} Брандмауэр Windows блокирует порты \texttt{pioneer\_sdk}: 8001 (команды MAVLink, UDP) и 8888 (видео с камеры).
    \item \textbf{Решение:} Разрешите Python доступ в частные сети.
\end{itemize}

\subsection{Ошибки в коде Python}

\subsubsection{IndentationError}
Python очень чувствителен к отступам. Не смешивайте табы и пробелы. Используйте 4 пробела для отступа.

\subsubsection{OpenCV Error: Assertion failed}
Обычно означает, что функция получила пустое изображение (None).
\begin{itemize}
    \item Проверьте, что \texttt{camera.get\_cv\_frame()} не вернул None.
    \item Проверьте, что файл изображения существует (если читаете с диска).
\end{itemize}

\subsubsection{ValueError: The truth value of an array \dots{} is ambiguous}
Детектор возвращает \texttt{ids} массивом numpy или \texttt{None}. Запись \verb|if ids:| для массива из нескольких элементов падает с этой ошибкой. Проверяйте явно: \verb|if ids is not None:|, а номер берите как \verb|int(ids[i][0])|.

\begin{taskbox}[Найдите три ошибки]
\emph{Симулятор.} Программа должна напечатать номера всех меток в кадре, но падает или печатает не то:
\begin{verbatim}
frame = camera.get_cv_frame()
corners, ids, _ = detector.detectMarkers(frame)
if ids:
    for i in range(len(ids) + 1):
        print("ID:", ids[i])
\end{verbatim}
Исправьте её и проверьте над меткой 12 и вдали от меток (у $Y = 5$, как в миссии Python №\,85).
\expected{над меткой печатается \texttt{ID: 12} (число, а не \texttt{[12]}), вдали от меток --- «меток нет», и программа не падает; кадр \texttt{None} тоже обработан до вызова детектора.}
\mission{Python}{81, 84--85}
\end{taskbox}

\subsection{Особенности симулятора}
Программа, которая работает в симуляторе, почти без изменений летит на настоящем дроне --- но есть отличия, о которых стоит помнить в обе стороны.
\begin{itemize}
    \item \textbf{\texttt{Camera(view="bottom")}} --- расширение симулятора. Настоящий \texttt{pioneer\_sdk} не знает параметра \texttt{view}: на дроне уберите его и подключайте камеру, которая установлена на аппарате.
    \item \textbf{Модель OpenCV.} В браузере \texttt{cv2} умеет получать и показывать кадр, рисовать и искать ArUco. \texttt{cv2.imwrite} ничего не сохраняет, \texttt{cv2.solvePnP} всегда возвращает \texttt{(False, None, None)}, а функций вроде \texttt{adaptiveThreshold}, \texttt{inRange} или калибровки нет --- обращение к ним даёт \texttt{AttributeError}. Такие программы запускайте на компьютере.
    \item \textbf{Словарь.} Симулятор сравнивает словарь детектора со словарём меток по имени. Метки миссий нарисованы из \texttt{DICT\_6X6\_250}; детектор с другим словарём (даже \texttt{DICT\_6X6\_50}) вернёт \texttt{ids = None}.
    \item \textbf{Метка «не видна».} Модель находит метку, только если она целиком в поле зрения камеры и ничем не закрыта. Метка у края кадра детектором не возвращается --- это нормальный исход (миссии №\,79 и 85).
    \item \textbf{Устаревший кадр.} Перед решением берите новый кадр: после \texttt{go\_to\_local\_point} дождитесь \texttt{point\_reached()} и дайте корпусу успокоиться (около 1,5~с), иначе кадр снят ещё в движении.
    \item \textbf{\texttt{set\_manual\_speed}} действует около 0,5~с: команду повторяют в цикле, а в конце явно отправляют нулевую скорость.
\end{itemize}

\subsection{Аварийные ситуации}

\subsubsection{Дрон не реагирует на Land}
Если программная посадка не срабатывает:
\begin{enumerate}
    \item Не паникуйте.
    \item Остановите программу и выключите моторы: тумблер взвода на пульте (раздел~\ref{sec:rc-control}) или кнопка Disarm в наземной станции (если есть).
    \item Только если моторы не выключаются: аккуратно поймайте дрон снизу за аккумулятор (берегите пальцы от пропеллеров!) и отключите аккумулятор. Не рассчитывайте, что моторы остановятся сами, если перевернуть дрон, --- проверьте, есть ли такая защита в прошивке вашего аппарата.
\end{enumerate}

\begin{warnbox}[Внимание!]
Никогда не пытайтесь схватить дрон сверху или сбоку — пропеллеры вращаются со скоростью свыше 10 000 об/мин и могут нанести серьезные травмы.
\end{warnbox}
